\section{Discusión}
\label{sec:discusion}

% ---------------------------------------------------------------
\subsection{La maquinaria transactiva escala; la representación era el límite}
\label{sec:disc-escala}

El resultado central de esta versión es una separación limpia de responsabilidades.
Al pasar de tres a ocho agentes, \emph{ningún} componente de la maquinaria transactiva
requirió cambios de mecanismo: los directorios hetero-asociativos, la lectura B1, el
rechazo por containment, la actualización de directorio de todos los agentes y la
coordinación punto a punto operan idénticos, y el directorio visual alcanza el balance
perfecto (3.000 de 3.000 bits) con cero falsos enrutamientos.
Lo que sí golpeó el escalado fue la \emph{representación}: con la binarización por signo
heredada de la versión anterior, la precisión temprana de 8 clases caía al 50\,\% y dos
clases colapsaban a 0\,\%.
La causa no estaba en las memorias sino en el alfabeto: $\mathrm{sign}(v)$ usaba 2 de
los 16 niveles disponibles, y ocho categorías con léxico solapado (tres frutas, tres
animales) no caben en ese margen.
La cuantización por magnitud sobre una escala global (Sección~\ref{sec:pipeline-semantico})
recuperó el rango de 85--88\,\% de precisión temprana \emph{sin tocar una línea de la
maquinaria asociativa}.
La Figura~\ref{fig:espacio-semantico} muestra el espacio semántico resultante: los
solapamientos entre dominios persisten (son propiedades del léxico), pero la
representación ya no los agrava.

\begin{figure}[H]
    \centering
    \includegraphics[width=\textwidth]{fig11_semantic_space.png}
    \caption{Análisis del espacio semántico de los labels de las ocho clases:
    proyección PCA 2D con centroides de dominio y tokens ambiguos, similitud coseno de
    tokens problemáticos y ocupación de los 16 niveles de la cuantización por magnitud.}
    \label{fig:espacio-semantico}
\end{figure}

% ---------------------------------------------------------------
\subsection{La normalización necesaria sigue siendo B1, también a ocho agentes}
\label{sec:disc-b1}

La lectura B1 del directorio transactivo no corrige un defecto: refleja cómo está
construido el sistema.
Cuando los agentes especialistas acumulan números desiguales de interacciones, la
puntuación cruda del directorio favorece de forma constante al agente más activo con
independencia del contenido semántico de la consulta.
Dividir entre el conteo normaliza la evidencia a \textit{evidencia por interacción}, que
es una medida justa para comparar agentes con historiales de exposición distintos.

El escalado refuerza esta lectura en dos direcciones.
Primero, la brecha cruda-B1 persiste con $K = 8$ (80.0\,\% frente a 93.2\,\%): el sesgo
de masa no era un artefacto del caso de tres agentes.
Segundo, la entropía del directorio es alta incluso en la condición cruda (2.955 de
3.000 bits): la \emph{estructura} del directorio está bien formada y balanceada; el
sesgo vive en la \emph{lectura}, exactamente donde B1 actúa.
Las condiciones que manipulan síntomas (B2, C, D, E) vuelven a quedar por debajo, y la
combinación G no supera a B1 sola: una vez corregida la causa, las correcciones
auxiliares son redundantes.

% ---------------------------------------------------------------
\subsection{La dominancia de un agente era un efecto conjunto: encoder y fuente de conocimiento}
\label{sec:disc-dominancia}

Las corridas iniciales de 8 clases mostraban al agente \texttt{apple} capturando
$\approx$20\,\% de las victorias maduras y a \texttt{pear} colapsando, y la explicación
disponible era léxica: \texttt{apple} posee en ConceptNet el vocabulario genérico de
fruta (\textit{fruit}, \textit{red}, \textit{green}, \textit{round}, \textit{tree}).
Dos intervenciones sobre el vocabulario (extracción bidireccional con masa asociativa y
desempate por centroide) movieron poco ese reparto.
El hallazgo llegó por una vía inesperada: un reentrenamiento determinista del encoder
(motivado por el benchmark de costo) produjo un espacio latente con mejor separación de
\texttt{pear} y \texttt{tomato} respecto a \texttt{apple}, y con él la participación de
\texttt{apple} cayó de 20.0\,\% a 13.0\,\% (ideal 12.5\,\%) y \texttt{pear} subió de
$\approx$52\,\% a 86\,\% \emph{sin cambiar el vocabulario}.

La conclusión corregida es que la dominancia era un efecto \emph{conjunto} de la
separabilidad del espacio latente y de la asimetría léxica de la fuente de
conocimiento, no de la segunda por sí sola.
Este hallazgo tiene una lectura metodológica: en un sistema bidireccional, un sesgo
observado en el hemisferio semántico puede tener su causa parcial en el hemisferio
visual, porque el llenado imagen-etiqueta acopla ambos espacios.

Del trabajo sobre el vocabulario queda un resultado con valor propio: enriquecer el
vocabulario de una clase ayuda \emph{solo si el conocimiento nuevo coincide con el
lenguaje de las consultas}.
\texttt{tomato} mejoró al adquirir los labels que sus consultas usan
(\textit{sauce}, \textit{salad}, \textit{vegetable}); \texttt{pear} empeoró al adquirir
variedades botánicas (\textit{bosc}, \textit{anjou}) que el banco apenas menciona,
porque el vocabulario extra solo diluye su activación media en los tokens compartidos.
El conocimiento del especialista debe estar en el idioma en que el grupo pregunta;
acumular experticia que nadie consulta no mejora al grupo.
Un intento relacionado que \emph{no} funcionó y conviene registrar: podar el vocabulario
por similitud coseno al nombre de la clase falla porque el nombre es polisémico en el
espacio de embeddings (\texttt{cup} $\approx$ \textit{World Cup} podaba \textit{drink} y
\textit{mug}).

% ---------------------------------------------------------------
\subsection{Lectura desde la teoría de memoria transactiva de Wegner}
\label{sec:disc-wegner}

Los resultados de las Secciones~\ref{sec:calibracion} y~\ref{sec:formacion} releen los
mecanismos de Wegner~\cite{wegner1987transactive} a una escala de grupo mayor:

\begin{itemize}
    \item \textbf{Un grupo más grande necesita más interacción para conocerse.}
    Con tres agentes, 10 consultas por dominio bastaban para cruzar el 90\,\% de
    precisión madura en $k \approx 13$ interacciones; con ocho agentes ese mismo volumen
    se queda en 82.5\,\%, y hace falta el banco completo ($\approx$50 por dominio) para
    alcanzar el 93.2\,\%.
    El costo de formación del directorio crece con el número de especialidades y con el
    solapamiento entre ellas, no solo con el número de miembros.

    \item \textbf{El orden de las experiencias afecta la velocidad, no el destino.}
    Todas las condiciones de presentación convergen exactamente al mismo directorio
    (mismos conteos, misma entropía).
    La estructura final del sistema transactivo depende de \emph{qué} se presenció, no
    de \emph{en qué orden}; el orden determina cuándo el directorio se vuelve utilizable.

    \item \textbf{La comunicación distintiva forma mejor directorio.}
    Las clases cuyas consultas usan léxico exclusivo (\texttt{cup}, \texttt{tomato},
    \texttt{car}) alcanzan 98--100\,\%; las que preguntan con léxico compartido
    (\texttt{pear}, \texttt{horse}) quedan en 82--86\,\%.
    En términos de Wegner: la asignación de información funciona en la medida en que
    quien pregunta aporta las señas del dominio.
    La palanca no es más vocabulario del especialista sino mejores señas del que
    consulta, como muestra el resultado de enriquecimiento
    (Sección~\ref{sec:disc-dominancia}).

    \item \textbf{El rechazo colectivo escala.}
    Con ocho especialistas y siete clases contrarias por agente, la especificidad se
    mantiene exacta: 0\,\% de falsos aceptados en contenido en todo el rango de $N$, y
    cero falsos enrutamientos visuales.
    El grupo sigue sin inventar referentes.
\end{itemize}

% ---------------------------------------------------------------
\subsection{Limitaciones}
\label{sec:limitaciones}

\paragraph{(a) Fuga de consultas fuera de dominio en el hemisferio textual.}
La sonda de rechazo (Sección~\ref{sec:proto-completo}) muestra que 2 de 12 consultas
fuera de dominio rutean hacia algún agente porque comparten tokens con labels legítimos.
Con la cuantización por magnitud y directorios densos, el containment textual pasa con
más facilidad que en la versión binaria, y la frontera semántica queda expuesta.
Una doble compuerta de contenido (verificar que el agente destino contenga la consulta,
no solo que el directorio la señale) es la mitigación diseñada pendiente de evaluación.

\paragraph{(b) Cobertura visual residual limitada por el containment conjunto.}
El 25\,\% de rechazo del enrutamiento visual de test no proviene de huecos por
coordenada (625 de 656 imágenes tienen soporte marginal completo) sino de la conjunción
AND sobre las 64 coordenadas simultáneas, por lo que aumentar $\xi$ no lo reduce
($\xi$ satura en 2--3).
La palanca es densificar el llenado con augmentación más rica, manteniendo el criterio
conservador que hoy garantiza cero falsos enrutamientos.

\paragraph{(c) Calidad de la reconstrucción visual.}
El decoder, entrenado con un término de reconstrucción MSE dentro de una pérdida que
incluye una cabeza auxiliar de clasificación, produce reconstrucciones suaves sin
detalle de alta frecuencia (cf.\ la reconstrucción e imaginación en EAM de Pineda y
Morales~\cite{pineda2023imagery}).
Es un efecto de la función de pérdida del autoencoder, no del sistema de memorias, y
puede abordarse de forma independiente.

\paragraph{(d) Aprendizaje limitado después de formar el directorio.}
Una vez llenas las memorias de dominio y formado el directorio, los agentes dejan de
registrar contenido nuevo: la fase madura enruta y recupera con lo aprendido.
El sistema aún no modela \emph{transactive encoding} ni la reorganización de memorias
de contenido durante la operación madura.

\paragraph{(e) Partición a nivel de imagen.}
La partición train/test se realizó a nivel de archivo de imagen, no a nivel de objeto
físico de ETH-80.
Una validación más estricta deberá usar particiones leave-object-out para medir
generalización a instancias completamente no vistas.

\paragraph{(f) Dependencia de la calidad de la fuente de conocimiento.}
El vocabulario de ConceptNet es asimétrico entre clases, y el resultado de
enriquecimiento muestra que no puede corregirse mecánicamente: el techo de las clases
de léxico compartido depende de la distintividad de las consultas del banco tanto como
del vocabulario del agente.
Fuentes de conocimiento adicionales (definiciones, corpora específicos de dominio) son
la extensión natural.
